\section{Limitations}
\label{sec:limitations}

\kunai is a DSL specialized for packet filtering, not a
general-purpose eBPF language: stateful filtering, deep packet
inspection, and aggregation are out of scope, and P4 defines only
header layout and parser transitions, not pipeline mechanisms such as
actions and tables.

\kunai addresses every field by name and offers no escape hatch for a raw
byte offset, as pcap-filter's \texttt{ip[14:2]} does. This
exclusion is deliberate: only a named, schema-resolved access lets the
generator bound each read into a form the verifier accepts, which an
arbitrary run-time offset does not. The raw byte-offset hatch is thus a
one-directional gap by design, not a shortfall.
Sub-byte fields such as a SYN-flag test are still written by name as
\texttt{where tcp.flags \& 0x02 != 0}, with the generator masking a
byte-aligned covering window.

Our verifier-acceptance result is empirical, not a formal guarantee:
we confirmed acceptance and matching on six kernels for F1--F10 and the
regression suite, but acceptance for an arbitrary \kunai expression on
an arbitrary kernel version is future work. Because a variable-length
walk is translated into bounded bytecode, the element count it can walk
has a compile-time bound, and a filter that uses \texttt{bpf\_loop} is
subject to a kernel-version lower bound.

The generated bytecode is only as correct as the protocol definition:
the p4c parse check verifies P4 syntax, not that field offsets or
parser transitions match the relevant RFC. The definitions also have
coverage limits: Geneve option TLVs, for example, are not yet exposed
as filter targets. Finally, \kunai's bytecode is a host-agnostic
evaluator over a packet window, and each host adapter supplies the
layout it sees; at tc, where the kernel moves the outer VLAN tag into
skb metadata, reading that tag value is future work
(\S\ref{sec:eval-accept}).
